\chapter{Machine Learning Models}
\label{ch:ml-models}

\section{Overview}

The platform deploys three machine learning models, each addressing a distinct intelligence requirement:

\begin{table}[H]
\centering
\caption{ML Model Registry}
\label{tab:model-registry}
\begin{tabular}{@{}llllll@{}}
\toprule
\textbf{Model} & \textbf{Algorithm} & \textbf{Task} & \textbf{Features} & \textbf{Output} \\
\midrule
SLA Risk & GradientBoostingRegressor & Regression & 9 & Score [0,1] \\
OSS Anomaly & IsolationForest & Anomaly Detection & 5 & Binary + score \\
BSS Revenue & IsolationForest & Anomaly Detection & 5 & Binary + score \\
\bottomrule
\end{tabular}
\end{table}

All models are packaged as scikit-learn \texttt{Pipeline} objects (StandardScaler + model) and serialized with \texttt{joblib} for deployment. The AI service dynamically reloads models when the underlying files change, supporting continuous retraining.

\section{SLA Risk Prediction Model}

\subsection{Problem Formulation}

Given a time window of aggregated network KPIs, predict the probability of an SLA breach in the next monitoring cycle. The output is a continuous score in $[0, 1]$ where higher values indicate greater breach risk.

\subsection{Feature Engineering}

Nine features are computed from raw OSS telemetry within each 2-minute window:

\begin{table}[H]
\centering
\caption{SLA Risk Model Features}
\label{tab:sla-features}
\begin{tabular}{@{}cll@{}}
\toprule
\textbf{\#} & \textbf{Feature} & \textbf{Description} \\
\midrule
1 & mean\_throughput\_mbps & Average throughput in window \\
2 & std\_throughput\_mbps & Throughput variability \\
3 & mean\_latency\_ms & Average round-trip latency \\
4 & std\_latency\_ms & Latency variability \\
5 & max\_latency\_ms & Peak latency observed \\
6 & mean\_packet\_loss\_pct & Average packet loss ratio \\
7 & max\_packet\_loss\_pct & Peak packet loss \\
8 & mean\_active\_users & Average concurrent users \\
9 & mean\_signal\_rsrp\_dbm & Average signal strength \\
\bottomrule
\end{tabular}
\end{table}

\subsection{Model Architecture}

\begin{equation}
\text{Pipeline: } X \xrightarrow{\text{StandardScaler}} X' \xrightarrow{\text{GBR}(200, d{=}4)} \hat{y} \xrightarrow{\text{clip}} [0, 1]
\end{equation}

\textbf{Hyperparameters:}
\begin{itemize}
    \item \texttt{n\_estimators}: 200 (with early stopping at \texttt{n\_iter\_no\_change}=15)
    \item \texttt{max\_depth}: 4 (captures feature interactions without memorizing noise)
    \item \texttt{learning\_rate}: 0.05 (conservative for stable convergence)
    \item \texttt{subsample}: 0.8 (stochastic gradient boosting for regularization)
\end{itemize}

\subsection{Training and Evaluation}

Training data consists of 3,000 synthetic time-windows generated with realistic telecom distributions. The target variable is computed as a weighted combination of latency and packet loss features with Gaussian noise:

\[
y = \text{clip}\left(\beta_1 \cdot \text{latency} + \beta_2 \cdot \text{packet\_loss} - \beta_3 \cdot \text{throughput} + \epsilon, \; 0, \; 1\right)
\]

\subsection{Explainability}

The GradientBoosting model provides built-in \texttt{feature\_importances\_} based on Gini importance. This is critical for CEM operations---operators need to know which KPIs drive risk. The top drivers are consistently latency-related features, validating the model's alignment with SLA definitions.

\section{OSS Network Anomaly Detection}

\subsection{Problem Formulation}

Given per-record network KPIs, identify anomalous data points that indicate network degradation, equipment failure, or unusual traffic patterns. This is an unsupervised task since labeled anomaly data is unavailable in production.

\subsection{Features}

Five per-record OSS metrics: throughput\_mbps, latency\_ms, packet\_loss\_pct, active\_users, signal\_rsrp\_dbm.

\subsection{IsolationForest Configuration}

\begin{itemize}
    \item \texttt{n\_estimators}: 150 (number of isolation trees)
    \item \texttt{contamination}: 0.05 (expected 5\% anomaly rate)
    \item Preprocessing: StandardScaler for feature normalization
\end{itemize}

IsolationForest returns:
\begin{itemize}
    \item \texttt{predict()}: +1 (normal) or -1 (anomaly)
    \item \texttt{decision\_function()}: Raw anomaly scores (more negative = more anomalous)
    \item Normalized scores: Mapped to $[0, 1]$ where 1 = most anomalous
\end{itemize}

\section{BSS Revenue Anomaly Detection}

\subsection{Anomaly Patterns}

The BSS model targets three specific fraud/anomaly patterns in the Tunisian telecom market:

\begin{table}[H]
\centering
\caption{BSS Anomaly Patterns}
\label{tab:bss-patterns}
\begin{tabular}{@{}lllll@{}}
\toprule
\textbf{Pattern} & \textbf{Revenue} & \textbf{Voice} & \textbf{SMS} & \textbf{Churn} \\
\midrule
SIM Box Fraud & 150--500 TND & 800--2000 min & 300--1000 & High \\
Dormant SIM & 0--1 TND & 0--1 min & 0--5 & High \\
SMS Spam & Normal & Normal & 300--1000 & High \\
\bottomrule
\end{tabular}
\end{table}

\subsection{Features}

Five BSS metrics: revenue\_tnd, data\_used\_gb, voice\_min, sms\_count, churn\_risk. Same IsolationForest configuration as the OSS model.

\section{Model Deployment Architecture}

Models are deployed through a shared Docker volume:

\begin{enumerate}
    \item Jupyter notebooks train and export models to \texttt{/app/models/}
    \item The \texttt{notebooks\_models} Docker volume persists artifacts
    \item AI service mounts the same volume as read-only
    \item Dynamic reloading: AI service checks file modification times every 30 seconds
    \item If a model file is newer than the cached version, it is reloaded automatically
\end{enumerate}

A bootstrap script ensures default models exist at startup, so the pipeline can run immediately even before notebooks are executed.

\section{Continuous Training Strategy}

The sliding-window approach enables model improvement:

\begin{enumerate}
    \item Pipeline generates new data every 120 seconds
    \item Data accumulates in PostgreSQL and MinIO
    \item Notebooks can be re-run at any time to retrain on the full dataset
    \item Retrained models are saved to the shared volume
    \item AI service automatically picks up new models within 30 seconds
\end{enumerate}

This architecture supports the transition from synthetic to real Tunisie Telecom data without any code changes---only the data generation step in the pipeline needs to be replaced with the real data ingestion connector.
